\section{The original length-two Witt Frobenius automorphism}
\label{sec:witt-two-frobenius}

Retain the actual ring $W_2(k)$ for a perfect
characteristic-$p$ field. The relative Frobenius
lift uses its actual coefficient Frobenius
automorphism, rather than a characteristic-$p$
power map on the mixed-characteristic Witt ring.

\begin{definition}[The original truncated Witt Frobenius]
\label{def:witt-two-frobenius}
\leanok
\lean{EnvelopeProperty.DeligneIllusie.WittTwoFrobenius.frobenius,EnvelopeProperty.DeligneIllusie.WittTwoFrobenius.frobenius_coeff}
The original full Witt Frobenius induces an actual
ring endomorphism $\sigma:W_2(k)\to W_2(k)$,
with $\sigma(x)_i=x_i^p$ on both actual Witt
coordinates.
\uses{def:witt-two-reduction}
\end{definition}
\begin{proof}
In characteristic $p$ for the coefficient field,
the actual full Witt Frobenius raises each Witt
coordinate to its $p$th power. Hence it sends a
full vector whose first two coordinates vanish to
another such vector. Its composite with original
truncation therefore vanishes on the kernel of
truncation. Factor this actual composite ring map
through the surjective original truncation map.
The resulting endomorphism is the asserted one.
Representing a truncated vector by its full vector
with later coordinates zero gives the displayed
formula on each of its two coordinates.
\end{proof}

\begin{lemma}[Compatibility with the original residue Frobenius]
\label{lem:witt-two-frobenius-reduction}
\leanok
\lean{EnvelopeProperty.DeligneIllusie.WittTwoFrobenius.reduction_frobenius}
The same actual ring map satisfies
$\rho(\sigma(x))=\rho(x)^p$ for every original
Witt vector $x$.
\uses{def:witt-two-frobenius,def:witt-two-reduction}
\end{lemma}
\begin{proof}
The original reduction is the degree-zero Witt
coordinate, so this is exactly the degree-zero
instance of the proved full coordinate formula.
\end{proof}

\begin{theorem}[The actual Witt Frobenius is an automorphism]
\label{thm:witt-two-frobenius-equivalence}
\leanok
\lean{EnvelopeProperty.DeligneIllusie.WittTwoFrobenius.frobenius_bijective,EnvelopeProperty.DeligneIllusie.WittTwoFrobenius.frobeniusEquiv}
The same original endomorphism is bijective and
hence an actual ring automorphism of $W_2(k)$.
\uses{def:witt-two-frobenius}
\end{theorem}
\begin{proof}
If two images agree, their two coordinates have
equal $p$th powers. Injectivity of Frobenius in
the original perfect field gives equality of both
original coordinates, so the vectors agree.
For any actual vector $x$, take the vector whose
coordinate $i$ is the unique Frobenius root of
$x_i$ in $k$. The proved coordinate formula shows
that its image is exactly $x$. This proves
surjectivity, and the actual ring endomorphism
with these two properties gives the claimed
original ring automorphism.
\end{proof}
